\subsection{Faces with one relation}
\label{subsec:general}

In this subsection $\widehat X$ and $F$ satisfy (S1)--(S3) of Section~\ref{subsec:setting}. The
argument uses only the charges of the light states. Its flat model is the rigid model of the
root.

\begin{lemma}\label{lem:flat-model}
Let $T^{k-1}$ act on $V=\HH^k$ with charge lattice $\Lambda$ as in (S2), and put
$m=\sum_i|a_i|$. Then $T^{k-1}$ acts freely on $\mu_0^{-1}(0)\smallsetminus\{0\}$, and
$V/\!/\!/T^{k-1}$ is $\CC^2/\ZZ_m$ with its flat orbifold metric; its coordinate ring is
$\CC[X,Y,s]/(XY-c\,s^m)$ for a nonzero constant $c$.
\end{lemma}

\begin{proof}
Write $\mu_i\in\mathbb R^3$ for the hyperk\"ahler moment map of the $i$-th summand $\HH$, so
that $|\mu_i|=\frac12|w_i|^2$ and $\mu_0=\sum_iq_i\otimes\mu_i$ on the charge vectors $q\in\Lambda$.
If $\mu_0(w)=0$, each of the three components of $(\mu_1,\dots,\mu_k)$ annihilates $\Lambda$,
hence is a multiple of $a$. If some $w_i=0$ then $\mu_i=0$, and since $a_i\neq0$ the multiples
vanish, so $w=0$. Every nonzero point of $\mu_0^{-1}(0)$ therefore has all coordinates nonzero.
Write $Q$ for a $(k-1)\times k$ matrix whose rows are a basis of $\Lambda$, and let
$\theta\in\mathbb R^{k-1}/\ZZ^{k-1}$ act on the $i$-th summand with weight $Q^{\mathsf T}\theta$.
The stabiliser of a point with all coordinates nonzero is $\{\theta:\ Q^{\mathsf
T}\theta\in\ZZ^k\}$, which is trivial because $\Lambda$ is saturated, that is because the Smith
invariants of $Q$ are all $1$.

In one complex structure write $w_i=(z_i,\tilde z_i)$, with $z_i$ of charge $q_i$ and $\tilde
z_i$ of charge $-q_i$, and put $u_i=z_i\tilde z_i$. A monomial in the $z_i,\tilde z_i$ is
invariant if and only if the difference of its exponent vectors lies in $\ZZ a$, so the
invariant ring of $\CC[z,\tilde z]$ is generated by $u_1,\dots,u_k$ and
\[
 X=\prod_{a_i>0}z_i^{a_i}\prod_{a_i<0}\tilde z_i^{-a_i},\qquad
 Y=\prod_{a_i>0}\tilde z_i^{a_i}\prod_{a_i<0}z_i^{-a_i},
\]
and every invariant monomial has a unique normal form $u^\gamma X^n$ or $u^\gamma Y^n$, so the
single relation $XY=\prod_iu_i^{|a_i|}$ generates all relations. The components of the complex
moment map are the invariant linear forms $\sum_iq_iu_i$, $q\in\Lambda$, so by the Reynolds
operator the ideal they generate meets the invariant ring in the ideal they generate there; on
their zero set $u=s\,a$ for a scalar $s$. The invariant ring of the complex level set is
therefore $\CC[X,Y,s]/(XY-c\,s^m)$, where $c=\prod_i\operatorname{sign}(a_i)^{|a_i|}|a_i|^{|a_i|}$,
a normal domain: the $A_{m-1}$ surface singularity. By the Kempf--Ness
theorem~\cite{KempfNess:1979}, $V/\!/\!/T^{k-1}$ is homeomorphic, and biholomorphic in the
chosen complex structure, to the spectrum of this ring.

Finally, $V/\!/\!/T^{k-1}$ is a hyperk\"ahler cone of real dimension four, smooth off its vertex.
A Ricci-flat cone over a three-manifold has an Einstein link, which in dimension three has
constant curvature, so the cone is flat, $\CC^2/\Gamma$ with $\Gamma\subset SU(2)$, and the
coordinate ring identifies $\Gamma=\ZZ_m$.
\end{proof}

\begin{theorem}\label{thm:general}
Let $\widehat X$ and $F$ satisfy (S1)--(S3), and assume Hypothesis~\ref{hyp:regular}. Let
$x=(\phi,b,0)$ be a root point with $\phi\in F$ and $b$ a smooth point of $N$. Then, as germs
of stratified spaces,
\begin{equation}
 \mathcal V_x\;\cong\;(\mathcal K_\phi\times N)\;\cup\;
 (F_\phi\times N\times\CC^2/\ZZ_m),\qquad m=\sum_i|a_i|,
 \label{eq:thm-general}
\end{equation}
the Coulomb branch and the Higgs branch, glued along $F_\phi\times N\times\{0\}$, where $N$
stands for its germ at $b$. The hypermultiplet factor of the Higgs branch off its root is a
smooth quaternionic K\"ahler manifold of quaternionic dimension $h^{2,1}(\widehat X)+2$, there is
no mixed branch, and the metric of the hypermultiplet factor is the flat orbifold metric of
$\CC^2/\ZZ_m$ times the flat metric on $T_bN$ up to relative corrections of order $r^2$, where
$r$ is the distance from the root.
\end{theorem}

\begin{proof}
By clause (1) of Hypothesis~\ref{hyp:regular} the Coulomb branch through $x$, on which the
charged hypermultiplets vanish, consists of supersymmetric Minkowski vacua. By (S3) the vector
kinetic matrix is nondegenerate at $\phi$, so by~\eqref{eq:vacua} the moment maps vanish there;
moreover, $\mu|_N$ is parallel (Lemma~\ref{lem:fixed}) and vanishes at the nearby Coulomb
points, hence on $N$. So $\mu(p)=0$ at the point $p\in M$ over $b$ with vanishing charged
hypermultiplets. Since no neutral hypermultiplet is charged, $N$ is the component through $p$ of
the fixed set of $T^{k-1}$ in $M$, which is the $N$ of Section~\ref{sec:quotients}. By Lemma~\ref{lem:isotropy}, $T^{k-1}$ acts on $T_pM$ through $Sp(n)$,
trivially on $T_pN$ and with charge lattice $\Lambda$ on $V=\HH^k$. By
Lemma~\ref{lem:flat-model} the action is free on $\mu_0^{-1}(0)\smallsetminus\{0\}$, so
Theorem~\ref{thm:normal-form} gives the hypermultiplet germ $N\times(V/\!/\!/T^{k-1})=N\times
\CC^2/\ZZ_m$, the part off $N$ being a smooth quaternionic K\"ahler manifold of quaternionic
dimension $(h^{2,1}(\widehat X)+1)+k-(k-1)$. The metric statement is
Proposition~\ref{prop:tangent-cone} together with the flatness in Lemma~\ref{lem:flat-model}.

The vector conditions are the last equation of~\eqref{eq:vacua}. Off the root of the Higgs
branch $T^{k-1}$ acts freely, so $h^Ik_I=0$ forces the image of $h$ in the Lie algebra of
$T^{k-1}$ to vanish, that is $h\cdot C_i=0$ for all $i$; since $F$ is open in the annihilator of
the $C_i$ by (S1), the vector moduli lie in $F_\phi$. On the Coulomb branch every
$h\in\mathcal K_\phi$ is allowed. By Theorem~\ref{thm:normal-form} the zero set near $p$ is
modelled on $N\times\mu_0^{-1}(0)$, whose points off $N$ have all charged coordinates nonzero,
so there is no mixed branch.
\end{proof}

When all $a_i=\pm1$, the case of $k$ nodes whose vanishing three-spheres are homologous up to
sign, $m=k$ and the transverse type is $\CC^2/\ZZ_k$. This is the configuration for which
Greene, Morrison and Vafa found $\CC^2/\ZZ_k$ with gravity decoupled~\cite[Sect.~3,
eqs.~(3.13)--(3.18)]{Greene:1996Confinement} and expected its qualitative features to persist in
the four-dimensional hypermultiplet space~\cite[Sect.~6]{Greene:1996Confinement}. The germ in
Theorem~\ref{thm:general} is the rigid model of the $k$ light hypermultiplets: gravity does not
remove the singular point or change its type.

\subsection{The four-node roots of \texorpdfstring{$X_9$}{X9}}
\label{subsec:four-node}

We call the roots over the face $F$ of~\eqref{eq:face-def}, where $\wideX_9$ contracts to the
threefold $X_F$ with four nodes, the \emph{four-node roots}.

\begin{theorem}\label{thm:four-node}
Assume Hypothesis~\ref{hyp:regular} for $\wideX_9$ and $F$. Let $x=(\phi,b,0)$ be a root point
with $\phi\in F$ and $b$ a smooth point of $N$. Then, as germs of stratified spaces,
\begin{equation}
 \mathcal V_x\;\cong\;(\mathcal K_\phi\times N)\;\cup\;
 (F_\phi\times N\times\CC^2/\ZZ_4),
 \label{eq:thm-four-node}
\end{equation}
glued along $F_\phi\times N\times\{0\}$, with $F_\phi$ of dimension $2=h^{1,1}(X_{9,t})-1$. The
hypermultiplet factor $N\times(\CC^2/\ZZ_4\smallsetminus\{0\})$ of the Higgs branch off its root is
a smooth quaternionic K\"ahler manifold of quaternionic dimension $124=h^{2,1}(X_{9,t})+1$. There
is no mixed branch.
\end{theorem}

\begin{proof}
We check (S1)--(S3) with $k=4$. By Proposition~\ref{prop:face}(1) the curves of zero area on
$F$ are the four disjoint $(-1,-1)$-curves $C_1,C_2,e_1,e_2$. Their classes span a lattice of
rank three with the single relation~\eqref{eq:face-relation}, so $a=(1,1,-1,1)$, and
$\operatorname{ann}(C_1,C_2,e_1,e_2)$ has dimension $6-3=3=\dim F$, so $F$ is open in it. The
charge lattice is saturated, since the charge matrix~\eqref{eq:app-Q} has Smith invariants
$(1,1,1)$. The kinetic matrix is positive definite on $F$ by Proposition~\ref{prop:face}(3).
Theorem~\ref{thm:general} applies with $m=4$ and $h^{2,1}(\wideX_9)=122$.
\end{proof}

\begin{lemma}\label{lem:completion}
In the situation of Theorem~\ref{thm:four-node}, the germ at $x$ of the metric completion of
$\mathcal H^{\mathrm{reg}}$ is $F_\phi\times(N\times\CC^2/\ZZ_4)$, and it coincides with the
closure of $\mathcal H^{\mathrm{reg}}$ in $\mathcal V_x$.
\end{lemma}

\begin{proof}
Near $x$, $\mathcal H^{\mathrm{reg}}=F_\phi\times(\mu^{-1}(0)\smallsetminus N)/T^3$, and by
clause (3) of Hypothesis~\ref{hyp:regular} its hypermultiplet metric is bi-Lipschitz equivalent
to the quotient metric. The link of the flat model is connected, so by
Lemma~\ref{lem:reduced-metric} the completion of $(\mu^{-1}(0)\smallsetminus N)/T^3$ near the
root is $(\mu^{-1}(0)\cap U)/T^3$, which by Theorem~\ref{thm:normal-form} is the germ of
$N\times\CC^2/\ZZ_4$. Bi-Lipschitz equivalent metrics have homeomorphic completions, and the
completion of a product is the product of the completions. The neighbourhood $U$ may be taken
relatively compact and intrinsic distances dominate ambient ones, so the completion maps
homeomorphically onto the closure of $\mathcal H^{\mathrm{reg}}$ in $\mathcal V_x$.
\end{proof}

\begin{proposition}\label{prop:x9-leading}
In the situation of Theorem~\ref{thm:four-node}, let $R_0$ be the leading curvature of the transverse
slice $\mathcal Y_p$ at the root, whose germ is $\CC^2/\ZZ_4$ (Section~\ref{subsec:leading}). Then
\begin{equation}
 R_0=4\nu R^{(1)}+\mathcal W,
 \label{eq:x9-leading}
\end{equation}
with $R^{(1)}$ as in Proposition~\ref{prop:universal} and $\mathcal W=c_\ell^{-4}\ell^*(R_{\mathrm{hk}}|_{\ell(\HH)})$
($c_\ell^2=4$) a $\ZZ_4$-invariant anti-self-dual Weyl tensor on $\CC^2$, where $R_{\mathrm{hk}}$ is the hyperk\"ahler
part of the curvature of $M$ at the root, whose restriction to $V=\HH^4$ is $T^3$-invariant. The $\ZZ_4$-invariant anti-self-dual Weyl
tensors form a real space of dimension $3$, and every element arises as $\mathcal W$ from some
$T^3$-invariant algebraic curvature tensor of hyperk\"ahler type on $V$.
\end{proposition}

\begin{proof}
The first term is Proposition~\ref{prop:universal}, and $\mathcal W$ is the restriction of $R_{\mathrm{hk}}$ to the
quaternionic line $\ell(\HH)$. Algebraic curvature tensors of hyperk\"ahler type on $\HH^n$
correspond to quartic polynomials in the complex coordinates, compatibly with the real structure given
by $I_2$~\cite{Alekseevsky:1968,Salamon:1982}; restriction to $\ell(\HH)$ corresponds to pulling back the quartic by the
complex-linear map underlying $\ell$, and invariance under $T^3$ to invariance of the quartic. The
$T^3$-invariant quartics in $(z_i,\tilde z_i)$ are spanned by the ten products $u_iu_j$ and the two
generators $X$, $Y$ of~\eqref{eq:app-ring}. With $\ell(\alpha,\beta)$ as in Section~\ref{subsec:leading},
\[
 u_i\mapsto\pm\alpha\beta,\qquad X\mapsto-\alpha^4,\qquad Y\mapsto\beta^4,\qquad u_iu_j\mapsto\pm\alpha^2\beta^2,
\]
so the image is spanned by $\alpha^4$, $\beta^4$ and $\alpha^2\beta^2$. These are exactly the quartics
invariant under the generator $(\alpha,\beta)\mapsto(i\alpha,-i\beta)$ of the stabiliser $\ZZ_4$ of
$\ell(\HH)$, which acts on the monomials $\alpha^{4-l}\beta^l$ by $i^{4-2l}$. The real structure induced on the quartics in $(\alpha,\beta)$ by the quaternionic structure $I_2$
exchanges $\alpha^4$ and $\beta^4$ up to conjugation and fixes $\alpha^2\beta^2$, so the real invariant tensors are
$c\,\alpha^4+\bar c\,\beta^4+t\,\alpha^2\beta^2$, $c\in\CC$, $t\in\mathbb R$, a space of real dimension three.
\end{proof}

The three directions have a direct interpretation. The products $u_iu_j$, invariant under the full torus
$U(1)^4$, give $\alpha^2\beta^2$, the direction invariant under the triholomorphic circle
$(\alpha,\beta)\mapsto(e^{i\theta}\alpha,e^{-i\theta}\beta)$; it is the Taub--NUT direction, spanned by the curvature at
the nut of the Taub--NUT space, whose Gibbons--Hawking potential~\cite{GibbonsHawking:1978} has a nonzero
constant term (checked in Appendix~\ref{app:data}). The invariants $X$ and $Y$ of Lemma~\ref{lem:flat-model} have degree
$m$; because $m=4$ here they are quartic and give $\alpha^4$ and $\beta^4$, the two directions that break this
circle. For a general face with one relation, $\Gamma=\ZZ_m$ acts on $\ell(\HH)$ by
$(\alpha,\beta)\mapsto(\zeta\alpha,\zeta^{-1}\beta)$ and hence on the quartics $\alpha^{4-l}\beta^l$ by $\zeta^{4-2l}$, so the Weyl part lies in
a real space of dimension $5$ for $m=2$, $3$ for $m=4$ and $1$ for all other $m$.

The proof uses only the quaternionic K\"ahler manifold $M$ at the root and the charges, so the same
statement holds in the setting of Theorem~\ref{thm:4d}, and for the metric of the type~IIA
hypermultiplet moduli space $\mathcal M_H(X_t)$ when the identification (W-IIA) of
Section~\ref{subsec:gmv} is an isometry. In that sense Proposition~\ref{prop:x9-leading} makes the
``non-universal'' correction of~\cite[Sect.~6]{Greene:1996Confinement} precise at leading order: a term
fixed by supergravity and a term in a three-dimensional space (a number specific to $m=4$), on which
$T^3$-invariance imposes no further linear constraint. Whether every such $\mathcal W$ arises from a quaternionic K\"ahler manifold with this isometric
torus action is not addressed.

\subsection{The del Pezzo facet}
\label{subsec:wall}

Through the smoothing of Proposition~\ref{prop:face-smoothing}, $H^2(X_{9,t};\mathbb R)$ is
identified with $\operatorname{ann}(C_1,C_2,e_1,e_2)$, spanned by $H_2,H_3,H_4$, with the same
cubic form~\eqref{eq:app-cubic}; we use the coordinates $y_2,y_3,y_4$ on both. The facet
$F_0=\{y_2=0\}$ of $\overline F$ is the locus where the whole surface $E$ collapses; there the light sector is the rank-one $E_2$ theory together with the two
nodes~\cite{W26global,W26quantum}. It is not a
boundary of the Higgs branch. By~\cite[eq.~(5.25)]{W26quantum} the Mori cone of $X_{9,t}$ lies
in the cone spanned by the degrees $(1,1,0)$, $(0,1,0)$, $(0,0,1)$ in the basis
$(H_2,H_3,H_4)$, so
\begin{equation}
 \mathcal K(X_{9,t})\supseteq\{y_2+y_3>0,\ y_3>0,\ y_4>0\},
 \label{eq:smooth-kahler}
\end{equation}
which contains the relative interior of $F_0$. The regular Higgs branch therefore continues
across the facet.

\begin{theorem}\label{thm:wall}
Assume Hypothesis~\ref{hyp:regular} for $\wideX_9$ and $F$. Let $y_0=y_2H_2+y_3H_3+y_4H_4$ lie in
the region~\eqref{eq:smooth-kahler}, normalised by $J^3=6$, and let $x$ be the limit point of the
regular Higgs vacua as the charged hypermultiplet expectation values, equivalently the smoothing
parameter $t$, tend to zero, with the vector moduli fixed at $y_0$ and the neutral
hypermultiplets fixed at a smooth point $b\in N$. Then the germ at $x$ of the metric completion of
$\mathcal H^{\mathrm{reg}}$ is
\begin{equation}
 \bigl(\text{germ of }\mathcal K(X_{9,t})\cap\{J^3=6\}\text{ at }y_0\bigr)\times
 \bigl(N\times\CC^2/\ZZ_4\bigr),
 \label{eq:thm-wall}
\end{equation}
whose vector factor is a full two-dimensional neighbourhood of $y_0$. In particular this holds
at every point $y_0$ of the relative interior of the del Pezzo facet $F_0$.
\end{theorem}

\begin{proof}
By Lemma~\ref{lem:product}, near $x$ the regular Higgs branch is the metric product of the
vector germ and a fixed manifold $\mathcal M_H^{\mathrm{reg}}$, and the metric completion of such
a product is the product of the completions. At a point of $F$, Lemma~\ref{lem:completion}
gives the completion as $F_\phi\times(N\times\CC^2/\ZZ_4)$, so the completion of
$\mathcal M_H^{\mathrm{reg}}$ at the root is $N\times\CC^2/\ZZ_4$. Since
$\mathcal M_H^{\mathrm{reg}}$ does not depend on the vector moduli, the same holds over $y_0$.
\end{proof}

Theorem~\ref{thm:wall} says nothing about the Coulomb branch at $x$, which at $F_0$ contains the
interacting $E_2$ point, and it describes the completion of the regular Higgs branch, not the
Higgs branch of the interacting theory at the root. The roots over the rest of the
region~\eqref{eq:smooth-kahler}, and the boundary $y_2=-y_3$, beyond which Theorem~\ref{thm:wall} does not apply, are
described in Section~\ref{subsec:other-roots}.
